\documentclass[11pt]{article}
\usepackage[margin=1in]{geometry}
\usepackage{amsmath,amssymb,amsthm,mathtools,booktabs}
\usepackage{xcolor}
\usepackage[colorlinks,linkcolor=blue!50!black]{hyperref}
\newtheorem{theorem}{Theorem}
\newtheorem{proposition}[theorem]{Proposition}
\newtheorem{lemma}[theorem]{Lemma}
\newtheorem{corollary}[theorem]{Corollary}
\theoremstyle{definition}
\newtheorem{definition}[theorem]{Definition}
\newtheorem{claim}{Claim}
\theoremstyle{remark}
\newtheorem{remark}[theorem]{Remark}
\newcommand{\LR}{\mathrm{LR}}
\newcommand{\Puz}{\mathrm{Puz}}

\title{Q585 at the Schur point: there is no deterministic YBE train\\
on Knutson--Tao puzzles to compare with the crystal commutor}
\author{Clio}
\date{12 November 2026 (prove session)}

\begin{document}
\maketitle

\begin{center}\small
\fbox{\begin{minipage}{0.9\textwidth}\small
\textbf{Provenance} (PROTOCOL \S2.3).
\textbf{Author:} Clio. \textbf{Date:} 12 Nov 2026.
\textbf{Repository:} \texttt{clio-vega/proofs}.
\textbf{Code commit:} \texttt{4641a0145b29}, directory \texttt{code-1112-Q585/}
(all numbers below were produced by those scripts at that commit).
\textbf{Registry:} \texttt{proofs/registry/q585-ybe-train-vs-commutor.json}.
\textbf{Not yet sent} (prove session: no email). The send is queued in \texttt{OUTGOING-UNSENT.md}.
\end{minipage}}
\end{center}

\begin{abstract}
Q585 asked whether the bijection $\LR(\lambda,\mu;\nu)\to\LR(\mu,\lambda;\nu)$ that a
Yang--Baxter ``train'' argument on Knutson--Tao puzzles would induce at the $0/1$ (Schur)
point equals the Henriques--Kamnitzer crystal commutor. The brief required one check
before any comparison: that such a train exists as a \emph{deterministic} map. It does
not, at least not as a train built from the puzzle model's own crossing. The crossing of
a line from the NW side with a line from the NE side factors through a single edge,
so it has rank $3$ out of $9$ (Proposition~\ref{prop:rank}). It annihilates exactly the
input that the equivariant rhombus would carry (Lemma~\ref{lem:kill}). The bundle
crossing it generates never has a single output, and it realises the swap
$|\lambda\rangle\otimes|\mu\rangle\mapsto|\mu\rangle\otimes|\lambda\rangle$ only on the
two constant words (Claim~\ref{cl:bundle}, $N\le5$). So the pre-registered falsifier
(``one $c\ge2$ case with $T\ne C$ and $T$ deterministic'') cannot fire. Q585 is
\emph{ill-posed at the instrument}. The prediction ``yes'' is neither confirmed nor
refuted.

The comparison side \emph{is} built and checked. The commutor $C$ is computed by two
instruments that share no code path: crystal operators, and RSK plus evacuation. They
agree on $1308/1308$ highest-weight elements. A lesioned $C$ differs on $1298$ of them.
A separating control (lex-rank matching) differs from $C$ on all $67/67$ cells with
$c\ge2$. A second plausible control (reverse-lex matching) agrees with $C$ on all
$67$, and would have been blind.
\end{abstract}

\section{What was asked, and the gate}

Let $A=B(\lambda)$, $B=B(\mu)$ be $\mathfrak{gl}_n$ crystals and $\xi$ the
Sch\"utzenberger involution. The Henriques--Kamnitzer commutor is
\[
  \sigma_{A,B}(b\otimes c)=\xi_{B\otimes A}\bigl(\xi_B(c)\otimes\xi_A(b)\bigr),
\]
and on highest-weight elements it gives a bijection $C:\LR(\lambda,\mu;\nu)\to\LR(\mu,\lambda;\nu)$.
(The formula is taken from \texttt{0707.2248}, read first-hand at l.855 in BROWSE 2026-11-11.)
Q585 asks whether $C$ equals a map $T$ ``induced by the YBE train'' on puzzles at the
$0/1$ point. In my reading index, Q585 was posed as \emph{the map a YBE train argument
would induce}, and no source in that index defines $T$. So the brief's \S2 gate is the
first question: does a deterministic $T$ exist?

A train argument moves a defect through a lattice model by local moves. Each move is an
identity of partition functions on a small region with fixed boundary. When every
weight is $0$ or $1$, each identity is an equality of \emph{counts}. It induces a
canonical map on configurations only if the counts are at most $1$ on both sides. A
count of $2$ on each side gives an equality of numbers and no canonical bijection.
That is the sense of ``deterministic'' below.

\section{The puzzle model, re-derived and calibrated}

Labels are $0,1,2$ (with $2$ standing for \texttt{10}). An up-triangle $\Delta$ has
(bottom, left$/$, right$\backslash$) in
$\{(0,0,0),(1,0,2),(2,1,0),(1,1,1),(0,2,1)\}$. A down-triangle $\nabla$ has
(top, right$/$, left$\backslash$) in the same five triples.

\begin{lemma}[calibration]\label{lem:calib}
With $\nu$ on the S side read left to right, $\lambda$ on the NW side read bottom to
top, and $\mu$ on the NE side read bottom to top, all encoded as $01$-words in
$n\times m$ boxes, the number of fillings equals $c^\nu_{\lambda\mu}$ in all $340$
cases with $n=m=3$ (all $\lambda,\mu,\nu\subseteq 3\times3$ with $|\nu|=|\lambda|+|\mu|$,
including $c^{321}_{21,21}=2$).
\end{lemma}

\noindent The enumerator (\texttt{puzzle.py}) is my own. The LR side counts
highest-weight tableau words (\texttt{commutor.py}), so the two sides share no code.
The orientation convention was \emph{fitted}: of the $16$ variants (three side
reversals times two readings of $\nabla$) exactly one scores $340/340$, and every other
scores at most $169/340$ (\texttt{calib.py}). I record this as a convention fixed by
data, not derived. Everything below uses that variant.

\paragraph{The model as lines.} Glue each $\nabla$ to the $\Delta$ directly above it
along their common horizontal edge. This gives a diamond $\Diamond$. Every $\nabla$ is
used once, and the only $\Delta$'s left over are the $N$ in the bottom row. A diamond is a
crossing of two lines: an \emph{SE-line} enters through its upper-left $/$ edge and
leaves through its lower-right $/$ edge, and an \emph{SW-line} enters upper-right
$\backslash$ and leaves lower-left $\backslash$. The SE-lines start on the NW side
(carrying $\lambda$) and the SW-lines start on the NE side (carrying $\mu$). Each pair
crosses at most once, and the bottom-row $\Delta$'s fuse one SE-line with one SW-line
into a letter of $\nu$. Exchanging $\lambda$ and $\mu$ means exchanging the two bundles,
so the only operator in the model that can carry information between them is the
diamond crossing.

\section{The crossing is not a braiding}

Write $D_0$ for the diamond, as the map
$(\text{SE}_{\rm in},\text{SW}_{\rm in})\mapsto(\text{SW}_{\rm out},\text{SE}_{\rm out})$
with $D_0[(\ell\ell,\ell r),(u\ell,ur)]=\#\{x:(x,u\ell,ur)\in\Delta,\ (x,\ell r,\ell\ell)\in\nabla\}$.

\begin{proposition}\label{prop:rank}
$D_0$ is a $0/1$ matrix of rank exactly $3$ on the $9$-dimensional space
$\mathbb{Q}^{\{0,1,2\}^2}$. In particular it is not invertible.
\end{proposition}
\begin{proof}
The upper bound is structural. A filling of the diamond is a choice of the middle
horizontal label $x\in\{0,1,2\}$ together with a $\Delta$ and a $\nabla$ on it, so
$D_0=\sum_x |\nabla_x\rangle\langle\Delta_x|$ factors through the $3$-dimensional label
space of that edge, and $\operatorname{rank}D_0\le3$. The entries are $0/1$ because the
$\Delta$ triple is determined by (left, right) whenever it exists: no two of the five
triples share both side labels. Equality: the inputs $(0,0)$, $(1,0)$ and $(0,2)$
force $x=0$, $x=2$ and $x=1$ respectively. Their output vectors are
$\{(0,0),(1,2)\}$, $\{(0,1)\}$ and $\{(1,1),(2,0)\}$, which have disjoint supports, so
they are linearly independent. (The full table is printed by \texttt{bundle.py}. SymPy
gives rank $3$, \texttt{local.py}.)
\end{proof}

\begin{lemma}\label{lem:kill}
$D_0$ annihilates the input $(\text{SE},\text{SW})=(0,1)$. It also annihilates $(1,2)$,
$(2,0)$ and $(2,2)$.
\end{lemma}
\begin{proof}
A filling needs $(x,u\ell,ur)\in\Delta$. No $\Delta$ triple has (left, right) equal to
$(0,1)$, $(1,2)$, $(2,0)$ or $(2,2)$. The side pairs that do occur are $(0,0)$,
$(0,2)$, $(1,0)$, $(1,1)$ and $(2,1)$.
\end{proof}

\begin{remark}
The input $(0,1)$ is where an SE-line carrying $0$ meets an SW-line carrying $1$. That
is exactly the configuration of the Knutson--Tao equivariant rhombus, whose weight
$y_j-y_i$ vanishes at the Schur point. This identification is from memory of the KT
equivariant rule and was not re-read this session. So the zero row in
Lemma~\ref{lem:kill} is the Schur-point shadow of the equivariant piece. The $0/1$
point and an invertible crossing are bought at the same moment, and cannot both be
had.
\end{remark}

\begin{claim}[computed]\label{cl:bundle}
Let $B_N$ be the bundle crossing in which each of $N$ SE-lines crosses each of $N$
SW-lines once by $D_0$, acting on pairs of $01$-words $(s,t)$ with equal numbers of
$1$'s. For $2\le N\le5$:
\begin{enumerate}
\item $B_N(s,t)$ is never a single configuration with multiplicity $1$. The number of
  distinct outputs reaches $5,12,49,152$, and $B_N$ kills $1/6$, $6/20$ and $28/70$ of
  the pairs for $N=2,3,4$.
\item Restricted to $01$-outputs, $B_N(s,t)$ equals exactly $(t,s)$ once only for the
  two constant pairs $s=t=0^N$ and $s=t=1^N$, where the swap is the identity.
\end{enumerate}
\end{claim}

\noindent So the natural train, which carries the $\lambda$-bundle across the
$\mu$-bundle with the model's own crossing, is not deterministic, and it does not
realise the exchange on any input that carries information. This is the content of
the gate.

\begin{corollary}[Q585 is ill-posed at the instrument]
Within the Schur-point Knutson--Tao model, no train built from its own crossing induces
a map $\LR(\lambda,\mu;\nu)\to\LR(\mu,\lambda;\nu)$. Hence the comparison ``$T=C$?''
has no left-hand side, and the pre-registered falsifier cannot fire.
\end{corollary}

\paragraph{What this does \emph{not} show.} It does not rule out a train using a defect
$R$ \emph{foreign} to the model, one that is not a product of puzzle pieces. A sweep
over all weight-conserving $0/1$ defects satisfying YBE and bootstrap relations with
$(D_0,\Delta)$ would be the next instrument. I did not run it, because I could not
write down which YBE such a defect must satisfy without a source that states the train
argument. That source is the reading task Q585 now reduces to (see \S\ref{sec:next}).

\section{The commutor side, built and controlled}

\paragraph{Instruments.} The tensor product $b\otimes c$ is the concatenated word
$\mathrm{read}(b)\,\mathrm{read}(c)$, where $\mathrm{read}$ reads rows bottom to top.
The crystal operators use the bracket rule ($i{+}1=\,$`(', $i=\,$`)'). Then
$\LR(\lambda,\mu;\nu)=\{b\in SSYT(\lambda):\mathrm{read}(b)\,\mathrm{read}(u_\mu)$ is
highest weight of weight $\nu\}$.
\begin{itemize}
\item[$C_1$:] $\xi$ on a connected component by crystal operators only. Walk $x$ up to
  the highest weight, recording $x=f_{i_1}\cdots f_{i_k}u$. Walk $u$ down to the lowest
  weight $\ell$. Then $\xi(x)=e_{i_1^*}\cdots e_{i_k^*}\ell$, with $i^*=n-i$.
\item[$C_2$:] no crystal operators. $\xi$ on tableaux is evacuation by jeu de taquin.
  $\xi$ on a word $w$ is the word with $P=\mathrm{evac}(P(w))$ and $Q=Q(w)$, by inverse
  RSK.
\end{itemize}
Preliminary checks: $C_1$'s $\xi$ equals evacuation on $405/405$ tableaux, and is
involutive. Crystal operators commute with RSK ($Q$ preserved, and $P$ moved by the
same operator) on $9000/9000$ random checks.

\begin{claim}[computed]\label{cl:C}
For $n=4$ and all $1\le|\lambda|,|\mu|\le5$ ($1241$ cells, $1308$ highest-weight
elements, $67$ cells with $c\ge2$):
$C_1=C_2$ on $1308/1308$. Every image has the form $c'\otimes u_\lambda$ with
$c'\in\LR(\mu,\lambda;\nu)$ ($1308/1308$), and $C_{\mu,\lambda}\circ C_{\lambda,\mu}=\mathrm{id}$
($1308/1308$).
\end{claim}

\paragraph{Controls.}
\begin{itemize}
\item \emph{Lesion} (drop the outer $\xi_{B\otimes A}$): differs from $C_1$ on
  $1298/1308$. So the agreement in Claim~\ref{cl:C} is not vacuous.
\item \emph{A non-control I nearly used:} ``apply $\xi$ to each factor, flip, and raise
  to highest weight'' \emph{equals} $C$ identically. $\xi_B(c)\otimes\xi_A(b)$ is the
  lowest-weight element of its component when $b\otimes c$ is highest weight, and $\xi$
  of a lowest-weight element is the highest-weight element. A cheap ``flip'' is the
  commutor in disguise.
\item \emph{Flip control $F$} (lex-rank matching of $\LR(\lambda,\mu;\nu)$ with
  $\LR(\mu,\lambda;\nu)$): differs from $C$ on $67/67$ cells with $c\ge2$, so the
  sample separates. Reverse-lex matching \emph{agrees} with $C$ on $67/67$. With
  $c=2$ the two matchings are the two bijections, so this records that $C$ is
  consistently the anti-lex matching in this range. It is a fact about the sample, not
  claimed beyond it. Had I chosen reverse-lex as the control, it would have read
  ``$F=C$ everywhere''.
\item For $\lambda=\mu$, $C$ has no fixed point in any of the $13$ cells with $c\ge2$.
\end{itemize}

\noindent A note on the brief's sample: ``$n=3$ ($\mathrm{Gr}(k,3)$)'' carries zero
bits, because every Littlewood--Richardson coefficient for $\mathrm{Gr}(k,3)$ is at most $1$.
The first case with $c=2$ needs $\nu=(3,2,1)$, so it lives in $\mathrm{Gr}(3,6)$ or
in $\mathfrak{gl}_3$ with $|\nu|=6$. The sample above contains it.

\section{What Q585 reduces to}\label{sec:next}
\begin{enumerate}
\item A \emph{source} for the train. Q585 asserted ``KTW prove $c^\nu_{\lambda\mu}=c^\nu_{\mu\lambda}$
  by moving an R-vertex through the puzzle''. I could not reconstruct that argument
  from the model, and the model's own crossing rules it out (above). The KTW work I can
  name that has a deterministic tropical move is the octahedron-recurrence proof of the
  LR rule. Henriques--Kamnitzer's octahedron paper (ID recalled as
  \texttt{math/0408114}, \emph{not in my index, unverified}) defines a commutor on
  hives by the tropical octahedron recurrence. If, as I recall, they prove it equals
  the crystal commutor, then the hive form of Q585 is a theorem in print, and the
  remaining question is transport through a puzzle--hive bijection. That is a reading
  task, graded \texttt{recalled}, and nothing in this note rests on it.
\item The $C$ instrument is ready for that comparison: two instruments, a lesion, and
  a separating control, over $67$ cells with $c\ge2$.
\end{enumerate}

\section{Verification summary}
\begin{tabular}{@{}lll@{}}\toprule
item & script & result\\\midrule
puzzle counts $=c^\nu_{\lambda\mu}$ & \texttt{calib.py} & $340/340$ (1 of 16 variants; others $\le169$)\\
$\operatorname{rank}D_0$ & \texttt{local.py} & $3$ (structural bound $\le3$)\\
$B_N$ single-output & \texttt{bundle.py} & $0$ pairs, $N=2..5$\\
$B_N$ clean swap & \texttt{bundle.py} & constant words only\\
$\xi$: crystal vs evacuation & \texttt{crystal.py} & $405/405$\\
RSK--crystal compatibility & \texttt{crystal.py} & $9000/9000$\\
$C_1=C_2$ & \texttt{commutor.py} & $1308/1308$\\
lesion differs & \texttt{controls.py} & $1298/1308$\\
$F_{\rm lex}\ne C$ & \texttt{controls.py} & $67/67$ cells\\\bottomrule
\end{tabular}

\end{document}
